\documentclass{article}
\usepackage{amsmath,amssymb,amsthm}
\usepackage{algorithm,algorithmic}
\usepackage{fullpage}
\usepackage{xcolor}
\newtheorem{theorem}{Theorem}
\newtheorem{lemma}{Lemma}
\newtheorem{assumption}{Assumption}
\newcommand{\E}{\mathbb{E}}
\newcommand{\R}{\mathbb{R}}

\begin{document}

\section*{Random Block Coordinate Descent (RBCD) for Non‑Convex Smooth Functions}

\section*{Mathematical Formulation}
Let $f:\R^{d}\rightarrow\R$ be differentiable (possibly non‑convex).  
Partition the coordinate vector $w\in\R^{d}$ into $m$ (disjoint) blocks
\[
w=(w^{(1)},w^{(2)},\dots ,w^{(m)}),\qquad 
w^{(i)}\in\R^{d_i},\;\;\sum_{i=1}^{m}d_i=d .
\]
Denote by $\nabla_i f(w)\in\R^{d_i}$ the partial gradient w.r.t. block $i$ and by $U_i\in\R^{d\times d_i}$ the embedding matrix that inserts a block vector into the full space (i.e. $U_i v$ has $v$ in block $i$ and zeros elsewhere).

At iteration $t$ a block index $i_t\in\{1,\dots,m\}$ is drawn independently
from a distribution $p$; we use the uniform distribution
\[
\mathbb{P}(i_t=i)=\frac1m,\qquad i=1,\dots,m .
\]

Given a stepsize $\eta>0$, the update reads
\begin{equation}
\label{eq:update}
w_{t+1}=w_t-U_{i_t}\Bigl(\eta\,\nabla_{i_t}f(w_t)\Bigr)
\;=\;
w_t-\eta\,U_{i_t}\nabla_{i_t}f(w_t) .
\end{equation}
All other blocks remain unchanged.

\section*{Pseudocode}
\begin{algorithm}[H]
\caption{Random Block Coordinate Descent (RBCD)}
\begin{algorithmic}[1]
\STATE {\bf Input:} initial point $w_0\in\R^{d}$, stepsize $\eta>0$, 
maximum iterations $T$.
\FOR{$t=0,\dots,T-1$}
    \STATE Sample $i_t\sim\text{Uniform}\{1,\dots,m\}$.
    \STATE Compute block gradient $g_t\gets\nabla_{i_t}f(w_t)$.
    \STATE Update only block $i_t$:
    \[
    w_{t+1}=w_t-\eta\,U_{i_t}g_t .
    \]
\ENDFOR
\STATE {\bf Output:} $\{w_t\}_{t=0}^{T}$ (or a random iterate $w_{\tau}$).
\end{algorithmic}
\end{algorithm}

\section*{Dry Run Example}
Consider the one‑dimensional case $d=1$ (thus a single block) with
\[
f(w)=(w-3)^2,\qquad w_0=0,\qquad \eta=0.1 .
\]
Since there is only one block, the algorithm reduces to ordinary gradient descent.

\begin{center}
\begin{tabular}{c|c|c|c}
$t$ & $w_t$ & $\nabla f(w_t)$ & $f(w_t)$ \\ \hline
0 & $0$ & $2(0-3)=-6$ & $9$ \\
1 & $0-0.1(-6)=0.6$ & $2(0.6-3)=-4.8$ & $(0.6-3)^2=5.76$ \\
2 & $0.6-0.1(-4.8)=1.08$ & $2(1.08-3)=-3.84$ & $(1.08-3)^2=3.6864$ \\
3 & $1.08-0.1(-3.84)=1.464$ & $2(1.464-3)=-3.072$ & $(1.464-3)^2=2.3713$ \\
\end{tabular}
\end{center}
The objective value decreases at each iteration, illustrating the mechanics of \eqref{eq:update} in the simplest setting.

\newpage
\section*{Assumptions}
\begin{assumption}[Blockwise $L_i$‑smoothness]\label{ass:smooth}
For each block $i\in\{1,\dots,m\}$ there exists $L_i>0$ such that for all 
$w\in\R^{d}$ and any $v\in\R^{d_i}$,
\[
f\bigl(w+U_i v\bigr)\le 
f(w)+\langle \nabla_i f(w),v\rangle+\frac{L_i}{2}\|v\|^{2}.
\]
\end{assumption}
\begin{assumption}[Lower boundedness]\label{ass:lb}
The objective is bounded below:
\[
f_{\inf}:=\inf_{w\in\R^{d}} f(w)>-\infty .
\]
\end{assumption}
\begin{assumption}[Stepsize]\label{ass:stepsize}
The stepsize satisfies
\[
0<\eta\le \frac{1}{\bar L},\qquad 
\bar L:=\frac{1}{m}\sum_{i=1}^{m}L_i .
\]
\end{assumption}

All expectations $\E[\cdot]$ are taken w.r.t. the random block selections
$\{i_t\}_{t\ge0}$ and any other randomness in the algorithm.

\section*{Main Convergence Theorem}
\begin{theorem}[Expected stationarity of RBCD]\label{thm:main}
Under Assumptions \ref{ass:smooth}–\ref{ass:stepsize}, let $\{w_t\}_{t\ge0}$ be the sequence generated by \eqref{eq:update}. Then for any integer $T\ge1$,
\begin{equation}
\label{eq:rate}
\frac{1}{T}\sum_{t=0}^{T-1}\E\!\bigl[\|\nabla f(w_t)\|^{2}\bigr]
\;\le\;
\frac{2m\bigl(f(w_{0})-f_{\inf}\bigr)}{\eta\,T}.
\end{equation}
Consequently, if we pick an iterate $w_{\tau}$ uniformly at random from 
$\{w_0,\dots,w_{T-1}\}$, then
\[
\E\!\bigl[\|\nabla f(w_{\tau})\|^{2}\bigr]\le
\frac{2m\bigl(f(w_{0})-f_{\inf}\bigr)}{\eta\,T}=O\!\Bigl(\frac{1}{T}\Bigr).
\]
\end{theorem}

\section*{Theoretical Properties}
\begin{itemize}
\item \textbf{Stability.} The condition $\eta\le 1/\bar L$ guarantees the
coefficient $1-\frac{\eta L_i}{2}$ in the descent lemma is non‑negative for every block, preventing divergence.
\item \textbf{Hyper‑parameter sensitivity.} The bound \eqref{eq:rate}
improves linearly with the stepsize $\eta$ up to the limit $1/\bar L$; larger $\eta$ yields a smaller right‑hand side but violates the descent condition.
\item \textbf{Comparison to full‑gradient GD.} Setting $m=1$ recovers the classic
$O(1/T)$ rate for (non‑convex) gradient descent. With $m>1$, the factor $m$ appears because only one block is updated per iteration; the same computational budget as a full gradient evaluation yields the same $O(1/T)$ rate when measured per coordinate update.
\item \textbf{Special case – identical block Lipschitz constants.} If $L_i\equiv L$, then $\bar L=L$ and the bound simplifies to
\[
\frac{1}{T}\sum_{t=0}^{T-1}\E\!\bigl[\|\nabla f(w_t)\|^{2}\bigr]
\le \frac{2m\bigl(f(w_{0})-f_{\inf}\bigr)}{\eta T},
\qquad \eta\le\frac{1}{L}.
\]
\end{itemize}

\section*{Discussion}
The theorem shows that even though only a random subset of coordinates is
touched at each iteration, the expected squared norm of the full gradient
decays at the optimal $O(1/T)$ rate for non‑convex smooth optimization.
The proof relies solely on the blockwise smoothness property and does not
require convexity or variance reduction. Extensions (e.g., importance
sampling, adaptive stepsizes) can be handled by modifying the descent
lemma accordingly.

\newpage
\section*{Preliminaries (Definitions and Lemmas)}
\begin{lemma}[Block descent lemma]\label{lem:descent}
Let $i\in\{1,\dots,m\}$ and $w\in\R^{d}$. For any stepsize $\eta>0$,
\[
f\bigl(w-U_i(\eta\nabla_i f(w))\bigr)
\le 
f(w)-\eta\bigl(1-\tfrac{\eta L_i}{2}\bigr)\,\|\nabla_i f(w)\|^{2}.
\]
\end{lemma}
\begin{proof}
Apply Assumption \ref{ass:smooth} with $v=-\eta\nabla_i f(w)$:
\[
\begin{aligned}
f\bigl(w-U_i(\eta\nabla_i f(w))\bigr)
&\le f(w)+\langle\nabla_i f(w),-\eta\nabla_i f(w)\rangle
   +\frac{L_i}{2}\|\eta\nabla_i f(w)\|^{2}\\
&= f(w)-\eta\|\nabla_i f(w)\|^{2}
   +\frac{\eta^{2}L_i}{2}\|\nabla_i f(w)\|^{2}\\
&= f(w)-\eta\Bigl(1-\frac{\eta L_i}{2}\Bigr)\|\nabla_i f(w)\|^{2}.
\end{aligned}
\]
\end{proof}

\section*{Main Proof (Step‑by‑Step Derivation)}
We prove Theorem \ref{thm:main}. Throughout the proof, expectations are
taken w.r.t. the random block $i_t$ conditioned on the history up to time $t$.
All non‑trivial steps are labelled.

\bigskip
\noindent\textbf{[Step 1]} Apply Lemma \ref{lem:descent} to the random update
\eqref{eq:update} at iteration $t$:
\[
f(w_{t+1})
\le 
f(w_t)-\eta\Bigl(1-\frac{\eta L_{i_t}}{2}\Bigr)
   \|\nabla_{i_t}f(w_t)\|^{2}.
\tag{1}
\]

\noindent\textbf{[Step 2]} Take conditional expectation w.r.t. $i_t$.
Because $i_t$ is uniform,
\[
\mathbb{P}(i_t=i)=\frac1m,\qquad
\E\bigl[\|\nabla_{i_t}f(w_t)\|^{2}\mid\mathcal{F}_t\bigr]
   =\frac1m\sum_{i=1}^{m}\|\nabla_i f(w_t)\|^{2}
   =\frac1m\|\nabla f(w_t)\|^{2},
\]
and similarly
\[
\E\!\left[L_{i_t}\,\|\nabla_{i_t}f(w_t)\|^{2}\mid\mathcal{F}_t\right]
   =\frac1m\sum_{i=1}^{m}L_i\|\nabla_i f(w_t)\|^{2}
   \le \bar L\;\frac1m\|\nabla f(w_t)\|^{2}.
\tag{2}
\]
Insert (2) into (1) and use the linearity of expectation:
\[
\begin{aligned}
\E\!\bigl[f(w_{t+1})\mid\mathcal{F}_t\bigr]
&\le f(w_t)-\eta\Bigl(1-\frac{\eta\bar L}{2}\Bigr)
   \frac1m\|\nabla f(w_t)\|^{2}.
\end{aligned}
\tag{3}
\]

\noindent\textbf{[Step 3]} Remove the conditioning by taking total expectation.
Denote $F_t:=\E[f(w_t)]$. From (3) we obtain the deterministic recursion
\[
F_{t+1}\le F_t-\frac{\eta}{m}\Bigl(1-\frac{\eta\bar L}{2}\Bigr)
   \E\!\bigl[\|\nabla f(w_t)\|^{2}\bigr].
\tag{4}
\]

\noindent\textbf{[Step 4]} Choose $\eta$ satisfying Assumption \ref{ass:stepsize},
i.e. $\eta\le 1/\bar L$. Then $1-\frac{\eta\bar L}{2}\ge\frac12$ and (4) simplifies to
\[
F_{t+1}\le F_t-\frac{\eta}{2m}
   \E\!\bigl[\|\nabla f(w_t)\|^{2}\bigr].
\tag{5}
\]

\noindent\textbf{[Step 5]} Sum inequality (5) from $t=0$ to $T-1$ and telescope:
\[
F_T\le F_0-\frac{\eta}{2m}
   \sum_{t=0}^{T-1}\E\!\bigl[\|\nabla f(w_t)\|^{2}\bigr].
\tag{6}
\]

\noindent\textbf{[Step 6]} Use the lower boundedness Assumption \ref{ass:lb},
$F_T\ge f_{\inf}$. Rearranging (6) yields
\[
\sum_{t=0}^{T-1}\E\!\bigl[\|\nabla f(w_t)\|^{2}\bigr]
\le
\frac{2m}{\eta}\bigl(F_0-f_{\inf}\bigr)
=
\frac{2m}{\eta}\bigl(f(w_0)-f_{\inf}\bigr).
\tag{7}
\]

\noindent\textbf{[Step 7]} Divide both sides of (7) by $T$ to obtain the
average bound \eqref{eq:rate}:
\[
\frac{1}{T}\sum_{t=0}^{T-1}\E\!\bigl[\|\nabla f(w_t)\|^{2}\bigr]
\le
\frac{2m\bigl(f(w_0)-f_{\inf}\bigr)}{\eta\,T}.
\]
\hfill $\square$

\section*{Rate Derivation (With Constants)}
The bound \eqref{eq:rate} can be written as
\[
\mathbb{E}\!\bigl[\|\nabla f(w_{\tau})\|^{2}\bigr]
\le
\underbrace{\frac{2m}{\eta}}_{\displaystyle C(\eta,m)}
\;\frac{f(w_0)-f_{\inf}}{T},
\]
where $C(\eta,m)$ is linear in the number of blocks $m$ and inversely
proportional to the stepsize. Choosing the largest admissible stepsize
$\eta=1/\bar L$ minimizes $C$:
\[
C_{\max}=\frac{2m\bar L}{1}=2m\bar L.
\]

\section*{Special Cases}
\begin{enumerate}
\item \textbf{Identical smoothness:} $L_i=L$ for all $i$.
    Then $\bar L=L$ and the optimal stepsize $\eta=1/L$ leads to
    \[
    \E\!\bigl[\|\nabla f(w_{\tau})\|^{2}\bigr]
    \le \frac{2mL\bigl(f(w_0)-f_{\inf}\bigr)}{T}.
    \]
\item \textbf{Single block ($m=1$).} RBCD reduces to standard gradient descent,
    and the bound recovers the classical $O(1/T)$ rate for non‑convex smooth
    optimization.
\item \textbf{Importance sampling.} If we replace uniform sampling by a
    distribution $p_i\propto L_i$, the coefficient $m$ in \eqref{eq:rate}
    can be replaced by $\sum_i p_i^{-1}$, potentially improving the bound.
\end{enumerate}

\end{document}